% HD-F1-T39. GPT-6 Astra Pro, /root/transport; 22 November 2026 UTC.
\section[HD-F1-T39: an effective even symplectic rank cutoff]{HD-F1-T39: a finite rank bound for even symplectic groups}
\label{sec:f1-symplectic-even-quadratic-rank-cutoff}

Put \(L_0=3\cdot10^{30}-3\), as in T36. Every original
remaining T2-OTHER source with almost simple projective
socle \(\operatorname{PSp}_{2r}(q)\), for even \(q\), satisfies
\[
 \boxed{\qquad r\le 4L_0^2+7<10^{62}.\qquad}
 \tag{F1T39.1}
\]
The retained field bound is \(q\le27\), so here
\(q\in\{2,4,8,16\}\). Together with T38 this gives a
finite range of natural ranks for the entire remaining
symplectic branch. No existence assertion is made for the
retained range. The proof constructs the genuine quadratic
modules, their lower-level subspaces and their top
projectors, including all characteristic-two scalar factors.

\paragraph{1. Original operators and a literal rank-two commutator.}
Retain the original \(G\subset GL_K(V)\),
\(m=\dim_KV\ge27\), scalar centre \(Z\), perfect layer
\(L\), and actual surjection
\(\beta:H=\operatorname{Sp}_{2r}(q)\longrightarrow L\)
from T36. Suppose \(r\ge4L_0^2+8\). Every irreducible
constituent of \(V_{\mathbf C}|_H\) then has minimal
level \(1\le k\le L_0\) for the actual generator
\(\tau_A\oplus\zeta_A\), where
\[
 \tau_A(x)=q^{\dim\ker_A(x-I)},\qquad
 \zeta_A(x)=(-q)^{\dim\ker_A(x-I)},\qquad
 A=\mathbf F_q^{2r}.
 \tag{F1T39.2}
\]
Both are genuine characters; explicit models are given below.

The finite-field Steinberg Theorem 30 applies to type
\(C_r\) in characteristic two as well. At the present
rank there is no exceptional graph automorphism and no
exceptional covering group. The actual universal-cover
square of T36 therefore gives, for every original \(g\),
\[
 \alpha_g(h)=a h^{2^e}a^{-1},\qquad
 a\in\operatorname{GSp}(A),\qquad
 \beta\alpha_g=c_g\beta.
 \tag{F1T39.3}
\]
A diagonal action is represented by such an \(a\);
indeed its multiplier is a square over this even field.
This uses the finite-field, all-root-system statement of
Steinberg's theorem, rather than extending an odd-field
formula by notation. The uniqueness in the universal
central-cover square follows because the quotient of two
lifts is a homomorphism from the perfect cover to its centre.

For \(g\notin Z\), \(\alpha_g\ne1\), by the trivial
centralizer of the simple socle in the actual almost simple
projective group. A semilinear map fixing every projective
line is scalar linear: coordinate lines make it diagonal,
lines \(\langle e_i+e_j\rangle\) equate its diagonal
entries, and lines \(\langle e_i+t e_j\rangle\) make its
field automorphism fix every \(t\in\mathbf F_q\).
Choose a line moved by \(a\sigma_e\), a nonzero vector
\(v\) on it, and \(w=a(v^{2^e})\). If \(\lambda\) is
the multiplier of \(a\), then the actual transvections give
\[
 t_{v,c}(x)=x+c\langle x,v\rangle v,
 \quad c\ne0,\qquad
 h=\alpha_g(t_{v,c})t_{v,c}^{-1}
     =t_{w,c'}t_{v,c}^{-1},\quad c'=c^{2^e}/\lambda.
 \tag{F1T39.4}
\]
Writing \(\ell_u(x)=\langle x,u\rangle\), the exact formula is
\[
 (h-I)x=-c\ell_v(x)v+
 c'\bigl(\ell_w(x)-c\ell_v(x)\ell_w(v)\bigr)w.
 \tag{F1T39.5}
\]
Both pairs of vectors and functionals are independent; the
coefficient matrix has determinant \(-cc'\ne0\). Hence
\[
 \operatorname{rank}_A(h-I)=2,
 \qquad \ker(h-I)=v^\perp\cap w^\perp,
 \qquad \beta(h)=[g,\beta(t_{v,c})].
 \tag{F1T39.6}
\]
The signs in (F1T39.5) are interpreted in the given field.
No original coefficient-space operator is rescaled.

\paragraph{2. Two genuine quadratic plane models, with their phases.}
Fix a nontrivial additive character
\(\psi:\mathbf F_q\to\{1,-1\}\).
Let \(U_+\) have quadratic form \(Q_+(x,z)=xz\), and
let \(U_-\) be \(\mathbf F_{q^2}\), as a two-dimensional
\(\mathbf F_q\)-space, with \(Q_-(x)=x^{q+1}\).
Their polar forms are nondegenerate. For a nondegenerate
even-dimensional quadratic space \(U\), write
\[
 \epsilon(U)=q^{-\dim U/2}\sum_{u\in U}\psi(Q(u)).
 \tag{F1T39.7}
\]
Direct summation gives \(\epsilon(U_+)=1\) and
\(\epsilon(U_-)=-1\): the norm has one zero preimage
and \(q+1\) preimages of every nonzero scalar.
The sum is multiplicative under orthogonal direct sums.

Write \(A=X\oplus X^*\), with \(\dim X=r\).
The partial Fourier basis
\[
 e_{x,z}=q^{-r/2}\sum_{y\in\mathbf F_q^r}
                 \psi(y^t z)\delta_{(x,y)}
 \tag{F1T39.8}
\]
identifies the actual permutation module \(\mathbf C[A]\)
with \(\mathbf C[U_+^r]\). The lower Siegel element
\((x,y)\mapsto(x,y+Bx)\), for symmetric \(B\), acts by
\(\psi(z^tBx)\); a Levi element acts by simultaneous
column change. A rank-one Weyl element swapping \(x_i,y_i\)
acts by the normalized Fourier transform for the polar form
of \(Q_+\). These identities follow immediately by
substitution into (F1T39.8).

For \(U_-\) we use the genuine representation of the
unitary Heisenberg semidirect product in
Imai--Tsushima, Lemma 2.2. Its restriction to a Hermitian
space of dimension \(N\) has character
\((-1)^N(-q)^{\dim\ker(g-I)}\). This is a named
published input, with its supplied proof referring to
G\'erardin, Theorems 3.3 and 4.9.2. Here \(N=2r\),
and the natural \(\operatorname{Sp}(A)\)-matrices preserve
on \(A\otimes\mathbf F_{q^2}\) the Hermitian form
\[
 h((u,v),(u',v'))=\sum_i(u_i^qv_i'+v_i^qu_i').
 \tag{F1T39.9}
\]
The Heisenberg coordinates are \(((u,v),a)\), with
\(a+a^q=h((u,v),(u,v))\). The actual coordinate map
\[
 z=a+\sum_i u_i^qv_i\in\mathbf F_q
 \tag{F1T39.10}
\]
changes their multiplication into
\[
 (u,v,z)(u',v',z')=
 (u+u',v+v',z+z'+\operatorname{Tr}_{q^2/q}\sum_i v_i^q u_i').
 \tag{F1T39.11}
\]
Its Schr\"odinger representation on \(\mathbf C[\mathbf F_{q^2}^r]\)
is \(\psi(z)X(u)Z(v)\), where
\(X(u)\delta_x=\delta_{x+u}\) and
\(Z(v)\delta_x=\psi(\operatorname{Tr}\sum_i v_i^qx_i)\delta_x\).
The uniqueness of this central-character representation
identifies the published genuine extension with this space.
Consequently its \(H\)-character is exactly \(\zeta_A\).

We determine its operators, so this identification retains
all scalar phases. For either plane, and subsequently for
any orthogonal sum \(U\), set
\[
 D_B\delta_F=\psi\!\left(
   \sum_iB_{ii}Q(F_i)+
   \sum_{i<j}B_{ij}B_U(F_i,F_j)\right)\delta_F,
 \tag{F1T39.12}
\]
where \(B_U\) is the polar form and \(F_i\) are the columns.
For the norm plane, (F1T39.10) shows directly that this
operator implements the actual action of the lower Siegel
element on (F1T39.11). The column-permutation Levi
operators implement its actual Levi action. They form a
representation of the Siegel parabolic, so their ratio with
the genuine operators is a one-dimensional character of
that parabolic.

On the symmetric-matrix radical this ratio is a
\(GL_r(q)\)-invariant additive character. The trace pairing
identifies its dual with quadratic forms on \(\mathbf F_q^r\).
There is no nonzero invariant quadratic form for \(r\ge3\).
Indeed, transitivity on nonzero vectors makes its value a
constant \(c\) there. For \(q>2\), scalar homogeneity
forces \(c=0\). For \(q=2\), the three basis vectors and
their pairwise sums give polar pairings equal to \(c\),
and their triple sum has value zero, again forcing \(c=0\).
Thus the ratio is trivial on this radical. On the Levi it
factors through determinant: the elementary matrices in
\(SL_r(q)\) are commutators
\([e_{ik}(a),e_{kj}(1)]=e_{ij}(a)\).
Conjugating a one-coordinate diagonal element by the
corresponding Weyl element inverts it. Its Fourier
intertwiner does the same on the displayed permutation
operator, so the ratio satisfies \(\chi(a)=\chi(a^{-1})\).
Since \(q-1\) is odd, this determinant character is trivial.

The actual rank-one Weyl operator is therefore \(c\mathcal F_U\),
where
\[
 \mathcal F_U\delta_x=q^{-\dim U/2}
              \sum_y\psi(B_U(x,y))\delta_y.
 \tag{F1T39.13}
\]
For a plane, \(\mathcal F_U^2=I\) and the group relation
\((w_i u_{ii}(1))^3=I\) determines \(c\).
Completing the quadratic expression in the finite sum gives
\((\mathcal F_U D_Q)^3=\epsilon(U)I\).
The two relations give \(c^2=1\), \(c^3\epsilon(U)=1\),
hence \(c=\epsilon(U)\). Thus the actual Weyl phase is
\(+1\) on \(U_+\) and \(-1\) on \(U_-\).

Every nondegenerate quadratic plane over this field is
isometric to one of these two planes. An isotropic vector
and a polar dual give the hyperbolic plane after replacing
the dual \(v\) by \(v+Q(v)e\). Explicitly, for an
anisotropic plane \(Q(x,y)=ax^2+bxy+cy^2\),
nondegeneracy gives \(b\ne0\), and anisotropy gives
\(a\ne0\). The polynomial
\(T^2+(b/a)T+c/a\) is irreducible. Choose a root
\(t\in\mathbf F_{q^2}\) and \(z_0\) with
\(z_0^{q+1}=a\). Then the actual linear isomorphism
\[
 (x,y)\longmapsto z_0(x+t y),\qquad
 N_{q^2/q}\bigl(z_0(x+t y)\bigr)=Q(x,y)
\]
identifies it with \(U_-\). Splitting off nondegenerate polar planes gives
an orthogonal decomposition in all even dimensions.
Consequently an actual tensor summand of
\((\tau_A\oplus\zeta_A)^{\otimes k}\) is
\[
 \mathcal R_U=\mathbf C[\operatorname{Hom}(X,U)],
 \qquad \dim U=2k,\qquad D=\dim\mathcal R_U=q^{2rk},
 \tag{F1T39.14}
\]
with operators (F1T39.12), column changes, and
\(\epsilon(U)\mathcal F_U\) in each column.
An isometry \(s:U\to U'\) gives the actual intertwiner
\(\delta_F\mapsto\delta_{sF}\). The full group
\(S=O(U,Q)\) therefore acts by these permutations and
commutes with \(H\). No quotient of a Clifford group,
or a formal character replacement, is used here.

\paragraph{3. Exact code embeddings and all lower-rank constituents.}
For a totally singular \(N\subset U\) of dimension \(j\),
put \(U_N=N^\perp/N\). Define
\[
 i_N\delta_{\bar F}=q^{-rj/2}
                  \sum_{\Delta:X\to N}\delta_{F+\Delta}.
 \tag{F1T39.15}
\]
This is a well-defined isometry with image \(C_N\), the
functions supported on \(\operatorname{Hom}(X,N^\perp)\)
and constant on its \(\operatorname{Hom}(X,N)\)-cosets.
It is an actual \(H\)-intertwiner. The quadratic phase
in (F1T39.12) descends to \(U_N\), and column changes
preserve the displayed sums. Finite Fourier orthogonality
checks the Weyl operators, including their normalization.
To check their scalar as well, successively pair each
isotropic basis vector of \(N\) with a polar dual and
make that dual isotropic by the correction just used.
This gives \(U\simeq U_+^j\perp U_N\), and hence
\(\epsilon(U)=\epsilon(U_N)\). In particular
\[
 \dim C_N=q^{2r(k-j)},
 \qquad C_N\simeq\mathcal R_{U_N},
 \tag{F1T39.16}
\]
and the last module is an actual tensor summand of level
\(k-j\), including the trivial module when \(j=k\).

The full Siegel weight of \(\delta_F\) is the quadratic
form \(Q\circ F\), by (F1T39.12) and the nondegenerate
additive trace pairing. Its polar rank is
\(\operatorname{rank}(F^*B_U)\), an even integer.
We prove directly that, for \(r\ge2k+2\),
\[
 \mathcal R_U^{<2k}=
   \sum_{\substack{\ell\subset U,\ \dim\ell=1\\Q|_\ell=0}}C_\ell,
 \qquad
 E:=\dim\mathcal R_U^{<2k}<q^{2k+2r(k-1)}.
 \tag{F1T39.17}
\]
The left side is the sum of irreducible \(H\)-submodules
whose largest polar weight rank is less than \(2k\).
The dimension bound follows from (F1T39.16) and the bound
\(\#\{\text{lines in }U\}<q^{2k}\), once the equality
is proved.

Let \(K\) be an irreducible submodule of largest polar
rank \(d<2k\), and let \(\Phi\) be a nonzero vector of
a full quadratic weight of polar rank \(d\). For a map
\(F\) in its support put
\[
 M=\operatorname{im}F,\quad l=\dim M,\quad
 R=M\cap M^\perp,\quad j=\dim R.
 \quad d=l-j,\quad l+j\le2k.
 \tag{F1T39.18}
\]
Choose \(X_1=\ker F\) and a complement \(X_2\).
The partial Fourier transform in \(X_1\) is an actual
operator of the corresponding symplectic subgroup.
Its slice with \(F|_{X_2}\) fixed has Fourier support
only in \(\operatorname{Hom}(X_1,R^\perp)\): a vector
not orthogonal to \(R\) raises the polar rank of the
space containing \(M\) from \(d\) by at least two.
Finite Fourier orthogonality therefore makes the original
slice invariant under \(\operatorname{Hom}(X_1,R)\).
Because \(\dim X_1>0\) and the full quadratic weight
is fixed, this also forces \(Q|_R=0\). Otherwise adding
on one kernel direction a vector \(v\in R\) with
\(Q(v)\ne0\) would have the same nonzero coefficient
but a different quadratic weight. This step retains the
quadratic refinement in characteristic two.

More generally, if \(\Delta:X\to R\) and
\(F|_{\ker\Delta}\) is onto \(M\), choose
\(X_2\subset\ker\Delta\); the same partial Fourier
argument gives \(\Phi(F+\Delta)=\Phi(F)\).
This compares all full-image lifts on one branch even
when \(q=2\). Explicitly, let \(F,F'\) map onto the
same \(M\), with \(F'-F\in\operatorname{Hom}(X,R)\),
and choose \(\Delta\) uniformly in that Hom space.
Failure of \(F|_{\ker\Delta}\) to be onto is equivalent
to a nonzero intersection of the row spaces of \(F\)
and \(\Delta\). For any fixed nonzero row \(w\),
\[
 \Pr(w\in\operatorname{row}\Delta)
       \le(q^j-1)q^{-r}.
 \tag{F1T39.19}
\]
This follows by taking the union over nonzero linear
functionals on \(R\), each of which makes its composite
with the uniform \(\Delta\) uniform in \(X^*\).
Applying the same bound to \(F'\) and the translated
map \(\Delta+F-F'\), the probability of either failure
is at most
\[
 2(q^l-1)(q^j-1)q^{-r}
       <2q^{l+j-r}\le2q^{-2}<1.
 \tag{F1T39.20}
\]
The common map \(F+\Delta\) thus compares their
coefficients by the two actual partial Fourier maps.

Put \(j_0=(2k-d)/2\). For every branch with
\(j=j_0\), one has \(M=R^\perp\). For each such
\(R\) and each surjective quotient map
\(\bar F:X\to U_R\) of the given full weight, let
\(c_{R,\bar F}\) be the common coefficient of its
full-image lifts. Subtract exactly
\[
 c_{R,\bar F}\,q^{rj_0/2}i_R(\delta_{\bar F})
   =c_{R,\bar F}\sum_{\Delta:X\to R}\delta_{F+\Delta}.
\]
Thus the basis coefficient removed is
\(c_{R,\bar F}\), with the normalization in (F1T39.15) retained.
Full-image lifts exist because \(r\ge2k+2\).
Distinct maximal branches do not interfere: a full map
onto \(R'^\perp\) lying in \(C_R\)'s support gives
\(R\subset R'\), hence \(R=R'\).
The remainder \(\Phi'\) still has the same full weight,
belongs to the \(H\)-stable space
\(K+\sum_{\dim R=j_0}C_R\) of polar rank at most \(d\),
and has no supported map with \(j=j_0\).

If \(\Phi'\ne0\), take a supported map \(F\) of
maximal ordinary map rank. With \(X_1=\ker F\), its
slice is supported on \(\operatorname{Hom}(X_1,R)\):
maximality first forces the added image into \(M\),
and equality of full weights then forces it into
\(M^\perp\). The invariance above makes the slice a
nonzero constant times that subspace's indicator.
Its partial Fourier transform is nonzero at every map
\(X_1\to R^\perp\). Now \(\dim X_1\ge2\), and
\(R^\perp/R\) contains the nondegenerate space
\(M/R\) and a nondegenerate orthogonal complement of
dimension
\(2k-2j-d\ge2\), since \(j<j_0\).
Choose two vectors in that complement with nonzero polar
pairing and use them as images of two independent kernel
directions. The resulting Fourier coefficient is nonzero
and has polar rank \(d+2\), a contradiction.
Thus \(\Phi'=0\), and the irreducible \(K\) lies in
the sum of these code modules. Each \(C_R\) lies in
\(C_\ell\) for any singular line \(\ell\subset R\).
Conversely every \(C_\ell\) has maximal polar rank at
most \(2k-2\). This proves (F1T39.17).

An irreducible constituent of minimal tensor level \(k\)
in (F1T39.14) must therefore have polar rank \(2k\).
Otherwise its actual orthogonal projection from the sum
in (F1T39.17) is nonzero on one \(C_\ell\).
Its averaged Hermitian adjoint embeds it in the actual
level-\(k-1\) module \(\mathcal R_{U_\ell}\), contrary
to minimality. This constructs the lower-level comparison
map rather than replacing its character formally.

\paragraph{4. The top multiplicity-one projector.}
Set \(s=|O(U)|<q^{4k^2}\), and decompose the genuine
commuting action as
\(\mathcal R_U=\bigoplus_\tau E_\tau\otimes D_\tau\),
where \(E_\tau\) is an irreducible \(O(U)\)-module of
dimension \(a\). Choose a surjective map
\(F_0:X\to U\), and its full quadratic weight
\(q_0=Q\circ F_0\). The maps of this weight form an
actual \(O(U)\)-torsor: their common polar kernel is
\(\ker F_0\), and each is uniquely \(tF_0\) with
\(t\in O(U)\). The stabilizer of \(q_0\) in the
Levi \(GL(X)\) surjects onto \(O(U)\) on the quotient
\(X/\ker F_0\). Its action on this torsor is the
commuting right-regular action. Thus the \(q_0\)-weight
space in \(D_\tau\) is the irreducible dual
\(E_\tau^*\), of dimension \(a\), under this stabilizer.

Every rank-\(2k\) constituent contains a weight in the
single \(GL(X)\)-orbit of \(q_0\), and therefore
contains \(q_0\). Decomposing \(D_\tau\) into actual
\(H\)-irreducibles, the preceding irreducibility shows
that exactly one top constituent \(\eta_\tau\) occurs,
with multiplicity one. To compare the multiplicity space
with the original module, fix a nonzero \(e\in E_\tau\).
The evaluation \(D_\tau=\operatorname{Hom}_S(E_\tau,\mathcal R_U)\to\mathcal R_U\),
\(f\mapsto f(e)\), is an injective \(H\)-map, since
\(S e\) spans the irreducible \(E_\tau\). All other
constituents form a module \(B_\tau\) whose image lies
in \(\mathcal R_U^{<2k}\), so \(\dim B_\tau\le E\). Distinct \(\tau\) have distinct
top constituents by their displayed stabilizer modules.
In particular every minimal-level constituent selected
above is one of the \(\eta_\tau\).

The exact averaging projector and its trace normalization are
\[
 p_\tau={a\over s}\sum_{t\in O(U)}\overline{\tau(t)}\,R_U(t),
 \qquad
 \operatorname{Tr}_{D_\tau}(x)
 ={1\over s}\sum_{t\in O(U)}\overline{\tau(t)}
                                  \operatorname{Tr}R_U(x,t).
 \tag{F1T39.21}
\]
Indeed the trace of \(p_\tau R_U(x)\) is
\(a\operatorname{Tr}_{D_\tau}(x)\); this explains the
factor \(1/s\), including multiplicities.

\paragraph{5. Uniform traces at the original commutator.}
The translation and multiplication operators on
\(\mathbf C[U^r]\) form an orthogonal basis of its full
endomorphism space, labelled by \(A\otimes U\).
The displayed genuine \(H\times O(U)\) operators
conjugate these labels by \(x\otimes t\), with unit
phases allowed. Taking the trace of this conjugation gives
\[
 |\operatorname{Tr}R_U(x,t)|^2
 \le q^{\dim\ker(x\otimes t-I)}.
 \tag{F1T39.22}
\]
Only this inequality is needed; the phases are not
asserted to vanish. For \(t\ne I\), its actual
permutation action gives
\(\operatorname{Tr}R_U(I,t)=q^{r\dim\ker(t-I)}\le Dq^{-r}\).
For the rank-two \(h\) of (F1T39.6), (F1T39.22) gives
\[
 |\operatorname{Tr}R_U(h,I)|\le Dq^{-2k},\qquad
 |\operatorname{Tr}R_U(h,t)|\le Dq^{-r+1}\quad(t\ne I).
 \tag{F1T39.23}
\]
For the second inequality restrict \(h\otimes t-I\)
to \(\ker(h-I)\otimes U\); its rank there is
\((2r-2)\operatorname{rank}(t-I)\ge2r-2\).

Put \(B=aD/s\). Equations (F1T39.17) and
(F1T39.21)--(F1T39.23), with \(|\tau(t)|\le a\), give
\[
\begin{split}
 \dim\eta_\tau&>
 B\left(1-q^{4k^2-r}-q^{4k^2+2k-2r}\right),\\
 |\operatorname{Tr}\eta_\tau(h)|&<
 B\left(q^{-2k}+q^{4k^2-r+1}+q^{4k^2+2k-2r}\right).
\end{split}
 \tag{F1T39.24}
\]
The subtracted lower-part bound is legitimate even for
\(a>1\), since \(\dim B_\tau\le E\) and
\(E/B<q^{4k^2+2k-2r}\).
For \(r\ge4k^2+8\), \(k\ge1\), and \(q\ge2\),
\[
 { |\operatorname{Tr}\eta_\tau(h)|\over\dim\eta_\tau}
 < {1/4+1/128+1/262144\over1-1/256-1/262144}
 = {67585\over261119}< {7\over27}.
 \tag{F1T39.25}
\]
The final comparison is the exact integer difference
\(7\cdot261119-27\cdot67585=3038>0\).
All actual constituents have \(k\le L_0\), so the
bound holds for their direct sum with every original
restriction multiplicity. Thus
\(|\operatorname{Tr}_V\beta(h)|<7m/27\).

Average a Hermitian form over the original finite group.
For a unitary operator \(T\),
\(\operatorname{rank}(T-I)\ge(m-\operatorname{Re}\operatorname{Tr}T)/2\).
The literal commutator in (F1T39.6) also satisfies
\(\operatorname{rank}(\beta(h)-I)\le2\operatorname{rank}(g-I)\),
by expanding \([g,t]-I=(gt-tg)g^{-1}t^{-1}\).
Consequently every original nonscalar inertia satisfies
\[
                   \operatorname{rank}(g-I)>5m/27.
 \tag{F1T39.26}
\]
A nonidentity scalar has rank \(m\), so satisfies the
same bound. The original product-one generating tuple
has \(b\ge12\) nonidentity entries and the exact budget
\(\sum_j\operatorname{rank}(g_j-I)=2m+6\).
It follows that \(2m+6>20m/9\), or \(m<27\), the
required contradiction. This proves (F1T39.1).

\paragraph{6. Original-source and framework scope.}
Every map used above is on the original finite coefficient
module: the universal-cover pullback, the actual
transvection commutator, the genuine quadratic tensor
identification, the code embeddings and their averaged
adjoints, and the literal central projectors (F1T39.21).
The source curve, its parabolic cohomology, the prescribed
primitive sixfold corner and its Hermitian cup pairing
remain their original objects and maps. Their comparison
is the identity, with polarization multiplier one.
In particular no whole ambient factor is substituted for
the target corner. All Hilbert/cycle, Poincar\'e,
variational and Adams/GRR maps already attached to that
corner retain their original domains and scalar factors.
This is an exclusion of original source tuples above the
explicit rank, and creates no additional detector family.

The operative new published input is the genuine
unitary Heisenberg extension of Imai and Tsushima,
\emph{European Journal of Mathematics} 9 (2023), article 31,
Lemma 2.2 and its stated G\'erardin dependencies,
\nolinkurl{doi:10.1007/s40879-023-00620-5}.
The characteristic-two code and top-projector statements
needed here are proved in full above. They are compared
separately with the inspected author version of
Montealegre-Mora--Gross's Clifford tensor-power work;
no unverified central quotient or additional duality
statement from that comparison is an input to this proof.
